\documentclass[class=article,border=10pt]{standalone}
\usepackage{verbatim}
\usepackage{textcomp}
\usepackage{tikz}
\usetikzlibrary{shapes,arrows}
\usepackage[americanresistors,americaninductors,RPvoltages]{circuitikz}
\usepackage{amsmath}

\begin{document}

\begin{circuitikz} [american, scale=1]
    \draw  
    (0,0) to [sV, l=$V_{gen}$] (0,3)
    (0,3) to [D*, l_=1N4001] (3,3)
    % Capacitor en paralelo
    (3,3) to [C, l=$100\,\mu\text{F}$] (3,0)
    % Línea superior hacia la resistencia
    (3,3) -- (5,3)
    % Resistencia de carga
    (5,3) to [R, l=$4.7\,k\Omega$] (5,0)    
    % Línea inferior de retorno (GND)
    (0,0) to [short,-] (5,0)    
    ;
    
    % Etiquetas de medición del osciloscopio
    \draw[<-,>=stealth,blue] (0,3.1) -- (0,4) [thick] node[above] {$X$};
    \draw[<-,>=stealth,blue] (5,3.1) -- (5,4) [thick] node[right] {$Y$};
    \draw[<-,>=stealth,blue] (5.1,0) -- (6.1,0) [thick] node[above] {$GND_{osc}$};    

\end{circuitikz}

\end{document}